%%%%%%%%%%%%%%%%%%%%%%%%%%%%%%%%%%%%%%%%%%%%%%%%%%%%%%%%%%%%%%%%%%%%%%%%%%%
% Englische Fassung des Uebungsblatts zu Lehreinheit 3 (Teil A und B).
% Vorlage: handouts/LE03.tex. Nur in handouts.tex mit \def\wuzsprache{en}.
% Der Kruskal-Text (Aufgabe 1) ist eigens auf Englisch geschrieben und wird
% mit "python3 tools/kruskal_check.py handouts/en/LE03.tex" geprueft.
% Wer ihn aendert, muss das Skript erneut laufen lassen!
%%%%%%%%%%%%%%%%%%%%%%%%%%%%%%%%%%%%%%%%%%%%%%%%%%%%%%%%%%%%%%%%%%%%%%%%%%%
\begin{uebungsblatt}{The Counting Path}

\aufgabe{The counting path}[in class, individual work, 8 min.]
Silently choose a number between 1 and 10. Count that many words from the
beginning of the text; this is your starting word. Count its letters and
move on exactly that many words. Repeat. Punctuation, including
apostrophes, does not count. If the next jump would take you beyond the end
of the text, your current word is your \textbf{final word}. Keep it to
yourself. \emph{(Text written for this exercise, fictitious story.)}

\medskip
\noindent\fbox{\parbox{\dimexpr\textwidth-2\fboxsep-2\fboxrule\relax}{%
\linespread{1.25}\selectfont\raggedright
\hyphenpenalty=10000 \exhyphenpenalty=10000 % keine Trennungen: jedes Wort ganz
% KRUSKAL-TEXT-BEGINN
When Hanna Berger first walked into the old library of her town, she was
looking for one single book. Her aunt had told her that it held a report
about the winter in which the lake froze over completely and people could
walk across to the other side with dry feet. Hanna wanted to know whether
this story was true. She asked the librarian, a quiet man with a grey
beard, where she should begin. He smiled and gave her no title, only a
piece of advice: \enquote{Start anywhere, but write down where each book
sends you.} So she began with a chronicle of the local fishermen. In it she
found a reference to old church registers that recorded every baptism and
every burial. In the margin of the church registers there was a note about
unusual cold, with a reference to the weather records of a teacher. The
teacher had measured the temperature every morning and written it into
small notebooks. These notebooks lay in the cellar, in a box labelled
\enquote{Miscellaneous}. Hanna worked for three weeks at a long table under
a high window. She read, compared and took notes, and sometimes she found
statements that contradicted each other. A merchant wrote that the ice had
been thin, a priest claimed that people had danced on it. Hanna learned
that every witness had his own reasons for claiming something. She sorted
the reports by date, by origin and by the question of who could have seen
what. In the end almost all traces, no matter where she had started, led
to the same place. It was a thin notebook of the teacher, which listed many
nights of hard frost for one February, and next to them a sentence in neat
handwriting: \enquote{Today I walked with my pupils across the ice to the
far shore.} Hanna put down her pen. She now knew that her aunt's story had
a true core, and she knew something else as well: whoever searches long and
carefully often arrives at the same source.
% KRUSKAL-TEXT-ENDE
}}

\medskip
\noindent My starting number: \rule{15mm}{0.3pt}\qquad
My final word: \rule{45mm}{0.3pt}\qquad
Number of jumps: \rule{15mm}{0.3pt}
\end{uebungsblatt}

\begin{uebungsblattb}{2}{Searching and Documenting}

\aufgabe{Collecting search terms}[in class, individual work, 10 min.]
Break your research question from Unit~2 down into two or three concepts.
For each, collect terms in English and in a second language, synonyms,
broader and narrower terms. Mark where truncation (\texttt{*}) is useful:
\texttt{learn*} finds \emph{learn}, \emph{learning}, \emph{learner}.

\smallskip
\noindent\textbf{My research question:}
\schreiblinien{1}
\begin{center}
\renewcommand{\arraystretch}{2.0}
\begin{tabular}{|p{0.22\textwidth}|p{0.22\textwidth}|p{0.22\textwidth}|p{0.22\textwidth}|}
\hline
 & \textbf{Concept 1} & \textbf{Concept 2} & \textbf{Concept 3} \\ \hline
Term (English) & & & \\ \hline
Synonyms & & & \\ \hline
second language & & & \\ \hline
Broader/narrower terms & & & \\ \hline
\end{tabular}
\end{center}

\aufgabe{Search string and counter-search}[in class, 5 min.]
\begin{enumerate}
  \item Build a search string from the table: synonyms with \texttt{OR},
    concepts with \texttt{AND}.
  \item Write a second search string that specifically finds work that
    could \emph{contradict} your expected answer (terms such as
    \texttt{no effect}, \texttt{null result}, \texttt{failed replication},
    \texttt{critique}, or the counter-hypothesis itself).
\end{enumerate}
\noindent Search string 1:
\schreiblinien{2}
\noindent Search string 2 (counter-search):
\schreiblinien{2}

\aufgabe{Search log}[start in class, homework, 10\,\% of the portfolio grade]
Carry out a documented search in at least two databases. One row per
search. Record search terms, databases, number of hits and selection
criteria precisely enough for another person to repeat your search.

\begin{center}
\small
\renewcommand{\arraystretch}{2.4}
\begin{tabular}{|p{0.07\textwidth}|p{0.12\textwidth}|p{0.24\textwidth}|p{0.11\textwidth}|p{0.07\textwidth}|p{0.07\textwidth}|p{0.09\textwidth}|}
\hline
\textbf{Date} & \textbf{Database} & \textbf{Search string} & \textbf{Filters} &
\textbf{Hits} & \textbf{Screen\-ed} & \textbf{Se\-lect\-ed} \\ \hline
 & & & & & & \\ \hline
 & & & & & & \\ \hline
 & & & & & & \\ \hline
 & & & & & & \\ \hline
 & & & & & & \\ \hline
\end{tabular}
\end{center}
\noindent\textbf{Inclusion and exclusion criteria} (e.\,g.\ language, time
period, publication type, population):
\schreiblinien{2}
\noindent\textbf{Snowballing:} Which two key texts did you follow backwards
and forwards? What new material did this turn up?
\schreiblinien{2}

\aufgabe{Sorting hits}[in class, 10 min.]
Assign your first five hits to a level. The levels are a working aid for
the first screening, not a quality judgement in the individual case; that
follows in Unit~4.

\begin{center}
\small
\begin{tabular}{cp{0.78\textwidth}}
\toprule
\textbf{Level} & \textbf{Publication type} \\
\midrule
A & systematic review, meta-analysis, peer-reviewed journal article \\
B & academic book, handbook chapter, peer-reviewed conference paper \\
C & grey literature: preprint, thesis, report by a public agency or foundation \\
D & trade press, science journalism, textbook knowledge without references \\
E & blog, video, advertising, source without identifiable authorship \\
\bottomrule
\end{tabular}
\end{center}

\aufgabe{Annotated bibliography}[homework, due by Unit~4]
Comment on five sources you have found, with two to three sentences each
on \emph{key message}, \emph{method} and \emph{relevance to your
question}. At least one source should come from the counter-search. Manage
the sources in a reference manager (Zotero, Citavi or JabRef).

\smallskip
\noindent\textbf{Example of an annotation:}

\noindent\begin{minipage}{\textwidth}\small
Nickerson, R.~S. (1998). Confirmation bias: A ubiquitous phenomenon in many
guises. \emph{Review of General Psychology, 2}(2), 175--220.
\url{https://doi.org/10.1037/1089-2680.2.2.175}

\emph{Key message:} People preferentially seek, interpret and remember
information that supports an existing belief; the effect appears in
science, medicine, law, politics and everyday life. \emph{Method:} review
of experimental and historical findings, no data collection of its own.
\emph{Relevance:} Explains why my search must also look for contradicting
findings.
\end{minipage}

\begin{lehrhinweis}[Expected answers]
Task~1: the final word is \textbf{\enquote{source}} (last word of the text)
for every starting number from 1 to 10; according to the checking script
\texttt{tools/kruskal\_check.py}, all ten paths merge at word 23
(\enquote{told}) at the latest, with 68 to 71 jumps depending on the start.
Anyone who lands elsewhere has miscounted; recounting together is a small
lesson in reproducibility. Tasks~2/3: at least two concepts, English terms
present, synonyms joined with OR inside the brackets, AND between concepts;
the counter-search contains the counter-hypothesis or terms for null
results and critique. Task~4: the log is repeatable if date, database,
exact string, filters and number of hits are given; at least two
databases, one of them a subject database, Google Scholar at most as a
supplement. Task~6: annotations separate key message, method and
relevance; no mere restatement of the abstract.
\end{lehrhinweis}
\end{uebungsblattb}
